% A readable account of the compact formal proof of Theorem 16.
% Build: lualatex proof-report.tex (twice)
\documentclass[11pt,a4paper]{article}
\usepackage{fontspec}
\setmonofont{DejaVu Sans Mono}[Scale=0.82]
\usepackage{amsmath,amssymb,amsthm}
\usepackage[margin=2.4cm]{geometry}
\usepackage{booktabs,longtable,array}
\usepackage{fancyvrb}
\usepackage{xcolor}
\usepackage{tikz}
\usepackage[colorlinks=true,linkcolor=blue!50!black,urlcolor=blue!50!black,citecolor=blue!50!black]{hyperref}
\usepackage{enumitem}
\setlist{itemsep=2pt,topsep=4pt}

\newcommand{\leanus}{\textunderscore\allowbreak}
\newcommand{\lean}[1]{{\ttfamily\let\_\leanus #1}}
\emergencystretch=3em
\newcommand{\GT}{G_{\mathrm{T}}}
\newcommand{\GF}{G_{\mathrm{F}}}
\newcommand{\GB}{G_{\mathrm{B}}}
\newcommand{\T}{\mathrm{T}}
\newcommand{\F}{\mathrm{F}}
\newcommand{\B}{\mathrm{B}}
\DefineVerbatimEnvironment{code}{Verbatim}{fontsize=\small,frame=leftline,framerule=0.6pt,rulecolor=\color{gray!50},xleftmargin=1em}
\theoremstyle{definition}
\newtheorem{fact}{Fact}

\title{The compact formal proof of Theorem~16\\[2pt]
\large What the Lean statements say, how they are proved, and what is trusted}
\author{Companion to\\ \emph{The Trident: A Minimum Symmetric Clause Gadget for 3SAT to 3-Coloring}\\
M.~Cho, H.~Kim, and I.~Shin}
\date{November 2026}

\begin{document}
\maketitle

\begin{abstract}
\noindent
This note accompanies the Lean~4 package
\lean{compact-proof-trident\_minimum\_unique}. The package keeps only the
declarations that the statements of \lean{atlas/Statement.lean} use, and adds
three statements that say what the engine functions in those statements mean.
The note explains the statements in plain mathematics, follows the proof of
the formal Theorem~16 (\lean{trident\_minimum\_unique}) from the top down,
lists the trusted base, and states what is \emph{not} formalized. A reader who
accepts the definitions explained in Section~\ref{sec:model} and the trust
assumptions of Section~\ref{sec:trust} can check, from this note and the
statement file alone, what the build of the package establishes.
The formal verification was written with Claude (Anthropic) and is under
evaluation by the authors; it is not part of the article, whose Theorem~16 is
proved by the argument of its Section~IV.
\end{abstract}

\tableofcontents

% ---------------------------------------------------------------------------
\section{The claim}\label{sec:claim}

The paper's Definition~1 calls a graph $G$ an \emph{$S_3$-symmetric clause
gadget} for the clause $C=(x_0\lor x_1\lor x_2)$ if it consists of three
pairwise non-adjacent literal vertices $x_0,x_1,x_2$, a palette triangle
$\GT,\GF,\GB$ with $\GB$ adjacent to every literal, and auxiliary vertices,
together with a group action $\rho:S_3\to\operatorname{Aut}(G)$ with
$\rho(\sigma)(x_i)=x_{\sigma(i)}$ that fixes each anchor. The gadget is
\emph{valid} if, with the anchors colored $\T,\F,\B$, a proper 3-coloring
extending a literal coloring in $\{\T,\F\}^3$ exists exactly for the
assignments that satisfy $C$.

\paragraph{Theorem~16 (paper).} Every valid $S_3$-symmetric clause gadget has
at least six auxiliary vertices, and with six it is the trident, up to a graph
isomorphism that fixes each anchor and each literal vertex.

\medskip
The package proves this statement for a formal candidate space
(Section~\ref{sec:model}) by combining two proved reductions with seven
compiled re-runs of the search and one kernel-checked re-run
(Section~\ref{sec:proof}). The passage from Definition~1 to the candidate
space rests on the paper's Lemma~6 (orbit structure); it is not formalized
(Section~\ref{sec:gaps}).

% ---------------------------------------------------------------------------
\section{The formal model}\label{sec:model}

The engine's definitions are in \lean{atlas/Atlas.lean},
\lean{atlas/Framework.lean}, \lean{atlas/Kits.lean} and
\lean{atlas/Problems/ThreeColoring.lean}; the two labelings of the
trident are in \lean{ProofsDecide.lean}. The definitions that only Statements~5--7
mention (\lean{SpecExt}, \lean{rowInst}, \lean{S3}, \lean{LegalUniverse}, and
a few others) sit in the proof files next to the lemmas about them.
Statements~1--4 use 56 declarations, 257 lines of source; all seven
statements use 76 declarations, 349 lines. The README lists their line
ranges.

\subsection{Permutations}
A permutation of $\{0,1,2\}$ is its value table $(s(0),s(1),s(2))$
(\lean{Perm}). The list \lean{perms} holds the six tables,
\lean{app s i} evaluates $s(i)$, \lean{compose s t} is $s\circ t$, and
\lean{parity s} is $0$ for the three even permutations and $1$ otherwise.

\subsection{Vertices, orbit types and the action}
A vertex is either a literal \lean{.lit i} ($i<3$) or an auxiliary vertex
\lean{.aux p j}, the $j$-th vertex of \emph{part} $p$. A \lean{Universe} is a
list \lean{parts} of part sizes. Its vertex list is the three literals
followed by the vertices of every part. A permutation $s$ acts by
\lean{Universe.act}:
\[
\begin{array}{lll}
x_i &\mapsto x_{s(i)}, & \\
(p,j) &\mapsto (p,j) & \text{if part $p$ has size 1 (fixed vertex)},\\
(p,j) &\mapsto (p,\ j\oplus \operatorname{parity}(s)) & \text{if size 2 (sign action)},\\
(p,j) &\mapsto (p,\ s(j)) & \text{if size 3 (natural action)},\\
(p,j) &\mapsto (p,\ \operatorname{index}(s\circ \pi_j)) & \text{otherwise (size 6: regular action)},
\end{array}
\]
where $\pi_j$ is the $j$-th entry of \lean{perms}. These are, up to
isomorphism, all transitive $S_3$-sets: the orbit of a point has size
$6/|\mathrm{Stab}|\in\{1,2,3,6\}$, and for each size the action is
determined by the stabilizer.

\paragraph{Statement~6 checks this action.}
For every list of part sizes in $\{1,2,3,6\}$, the action of the candidate
universe fixes every vertex under the identity, satisfies
$s\cdot(t\cdot v)=(s\circ t)\cdot v$, maps the vertex list onto itself, sends
$x_i$ to $x_{s(i)}$, and fixes the three anchors
(\lean{candidate\_action\_is\_s3\_action}). \lean{Meaning.attach\_mulAction}
packages it as a Mathlib \lean{MulAction} of the group \lean{S3}, and
\lean{Meaning.attach\_edgeOrbit\_length} shows by the orbit--stabilizer
theorem that every edge orbit has size 1, 2, 3 or 6.

\subsection{The candidate space}
\lean{threeColoring.attach parts} returns the universe with parts
\lean{parts ++ [1, 1, 1]}: the three appended singleton parts are the anchors
$\GT,\GF,\GB$ (vertices \lean{.aux n 0}, \lean{.aux (n+1) 0},
\lean{.aux (n+2) 0} for $n=$ \lean{parts.length}). It also returns the
\emph{forced edges}: the palette triangle and the three edges from $\GB$ to
the literals.

\lean{Universe.edges} lists every unordered pair of distinct vertices except
literal--literal pairs, each in a normal form (\lean{normEdge}). This
includes pairs between a literal and $\GT$ or $\GF$, so the candidate space
is larger than the paper's standing assumption after Definition~1 requires.
\lean{actEdge s e} applies $s$ to both ends and normalizes.

A \emph{symmetric candidate} of composition \lean{parts} is an edge list
\lean{g} with
\begin{code}
abbrev SymCandidate (parts : List Nat) (g : List Edge) : Prop :=
  (∀ e ∈ g, e ∈ (threeColoring.attach parts).1.edges) ∧
  (∀ s ∈ perms, ∀ e ∈ g, (threeColoring.attach parts).1.actEdge s e ∈ g) ∧
  (∀ f ∈ (threeColoring.attach parts).2, normEdge f ∈ g)
\end{code}
that is, a graph on the universe with no literal--literal edge, invariant
under the action, and containing the palette triangle and the $\GB$ fan.

\subsection{Orbit compositions}
\lean{compositions k [1, 2, 3, 6]} enumerates the lists of part sizes of $k$
auxiliary vertices. \textbf{Statement~5} (\lean{compositions\_are\_all\_orbit\_compositions})
says that it lists exactly the non-decreasing lists over $\{1,2,3,6\}$ that
sum to $k$. Nothing is missing, so ``for every composition in the list''
means ``for every orbit composition''. The order of the parts is no
restriction, since parts are labels.

\subsection{The judge}
An \emph{input} is a number $\mathit{env}<8$ whose bit $i$ says whether
$x_i$ is true. Input $\mathit{env}$ satisfies $C$ iff $\mathit{env}\neq 0$.
For an input, \lean{rowInst parts g env} pins the anchors to the colors
$0,1,2$ (read $\T,\F,\B$), the true literals to $0$ and the false literals to
$1$; \lean{colorsOk} runs the backtracking 3-coloring test \lean{colorable}
on it.

\textbf{Statement~7} (\lean{gadgetK\_is\_validity}) says
\begin{code}
gadgetK threeColoring parts g false = some b ↔
  b = 1 ∧ ∀ env, env < 8 → (colorsOk (rowInst parts g env) = true ↔ env ≠ 0)
\end{code}
and \lean{Meaning.colorsOk\_rowInst\_iff} says that
\lean{colorsOk (rowInst parts g env) = true} holds iff there is a coloring
$f$ of the vertices with values in $\{0,1,2\}$ that agrees with every pin and
gives adjacent vertices different colors (the predicate \lean{SpecExt}).
Together: \lean{gadgetK threeColoring parts g false = some 1} holds exactly
when \lean{g} is a valid clause gadget in the sense of Definition~1, judged
on all eight inputs. The argument \lean{false} selects these eight rows;
with \lean{true} the judge uses only the four inputs $0,1,3,7$ (one per
number of true literals), which Section~\ref{sec:symred} shows to be
equivalent on symmetric candidates.

\subsection{The trident}
With \lean{parts = [3, 3]} the vertices are the literals, the parts
$(0,j)$ and $(1,j)$, and the anchors $\GT=(2,0)$, $\GF=(3,0)$, $\GB=(4,0)$.
\lean{tridentA} (\lean{ProofsDecide.lean}) is the trident with branch
vertices $p_j=(0,j)$ and upper triangle $u_j=(1,j)$: the triangle on part 1,
the edges $p_j u_j$, $x_j p_j$ and $p_j\GT$, the $\GB$ fan and the palette
triangle, 18 edges. \lean{tridentB} is the same graph with the roles of the
two parts exchanged.

% ---------------------------------------------------------------------------
\section{The seven statements}\label{sec:statements}

\lean{atlas/Statement.lean} restates the results; \lean{lake build} checks
them and prints their axioms.
\begin{enumerate}
\item \lean{trident\_is\_symmetric\_candidate}: \lean{SymCandidate [3, 3] tridentA}.
\item \lean{trident\_is\_valid}: \lean{gadgetK threeColoring [3, 3] tridentA false = some 1}.
\item \lean{no\_valid\_symmetric\_gadget\_below\_12\_vertices}: for every
$k\le5$, every composition of $k$ and every symmetric candidate \lean{g},
\lean{gadgetK threeColoring parts g false = none}.
\item \lean{valid\_symmetric\_gadget\_up\_to\_12\_vertices\_is\_trident}: for
every $k\le6$, composition and symmetric candidate \lean{g}, if
\lean{gadgetK} returns \lean{some b}, then \lean{parts = [3, 3]}, $b=1$, and
\lean{g} has the edge set of \lean{tridentA} or of \lean{tridentB}.
\item \lean{compositions\_are\_all\_orbit\_compositions} (Section~2.4).
\item \lean{candidate\_action\_is\_s3\_action} (Section~2.2).
\item \lean{gadgetK\_is\_validity} (Section~2.5).
\end{enumerate}
Read with Statements~5--7, Statements~3 and~4 say: among the $S_3$-symmetric
candidates with at most six auxiliary vertices (at most twelve vertices in
all), the only valid clause gadgets are the two labelings of the trident;
Statements~1 and~2 say that the trident is one. This is Theorem~16 in the
formal model. Statement~3 follows from Statement~4 (the composition
\lean{[3, 3]} sums to six), and Statement~4 is \lean{trident\_minimum\_unique}.

% ---------------------------------------------------------------------------
\section{The proof of \lean{trident\_minimum\_unique}}\label{sec:proof}

Every declaration written in the source is used by one of the seven
statements; the dependency check of Section~\ref{sec:checks} confirms it.
The proof of \lean{trident\_minimum\_unique} uses 278 declarations. It has
the following backbone: every dependency path from
\lean{trident\_minimum\_unique} to a declaration below a node passes through
that node (the dominator tree of its dependency graph; the numbers count the
declarations each node dominates, itself included).

\begin{code}
trident_minimum_unique (278)                   ProofsTheorem16
├─ gadgetK_s3_symmetry_reduction_threeColoring (51)  ProofsSym   §4.3
│    └─ colorsOk_row_act (38)
│         └─ colorable_rename (7)
├─ gadgetK_none_of_small_orbits (12)           ProofsSmallOrbits §4.1
├─ exists_graph_of_invariant (12)              ProofsCover       §4.4
├─ colorable_iff_spec (11)                     ProofsColor       §5
├─ threeColoring_kill (8)                      ProofsWalk        §4.2
│    └─ colorsOk_antitone (6)                  ProofsPrune
└─ trident_unique_threeColoring_invariant (5)  ProofsCover       §4.5
\end{code}

\paragraph{Case split.} \lean{trident\_minimum\_unique} fixes $k\le6$ and
rewrites \lean{compositions k [1, 2, 3, 6]} into its explicit value for each
$k$ (by \lean{simp}); the 24 compositions are then closed one by one:

\begin{center}
\begin{tabular}{@{}p{0.34\textwidth}p{0.3\textwidth}p{0.28\textwidth}@{}}
\toprule
compositions & argument & re-run\\
\midrule
sizes 1 and 2 only (16 compositions) & \lean{gadgetK\_none\_of\_small\_orbits} & none\\
\lean{[3]} & empty sweep & \lean{sym\_empty\_3} (kernel)\\
\lean{[1,3]}, \lean{[2,3]}, \lean{[1,1,3]}, \lean{[6]}, \lean{[1,2,3]}, \lean{[1,1,1,3]} & empty sweep & \lean{sym\_empty\_*} (compiled)\\
\lean{[3,3]} & the trident slice & \lean{trident\_slice\_found} (compiled)\\
\bottomrule
\end{tabular}
\end{center}

\subsection{Compositions of sizes 1 and 2}\label{sec:small}
\lean{gadgetK\_none\_of\_small\_orbits} (\lean{ProofsSmallOrbits.lean})
proves, for any number of auxiliary vertices, that no symmetric candidate
whose parts all have size 1 or 2 is valid. Let $c=(1,2,0)$, the 3-cycle
$x_0\mapsto x_1\mapsto x_2\mapsto x_0$. It is even, so it fixes every vertex
of a part of size 1 or 2 and every anchor (\lean{act\_cyc\_aux}). If $w$ is
adjacent to $x_2$, then $w$ is not a literal, so $c$ fixes $w$, and invariance
maps the edge $x_2w$ to $x_0w$: every neighbor of $x_2$ is a neighbor of $x_0$
(\lean{adj\_lit\_two\_to\_zero}). Take a proper coloring for the satisfying
input $(\F,\F,\T)$ (bit $2$ set, $\mathit{env}=4$) and recolor $x_2$ with
$\F$. Each neighbor of $x_2$ is adjacent to $x_0$, which has color $\F$, so
it avoids $\F$, and the result is a proper coloring for $(\F,\F,\F)$
(\lean{colorsOk\_row\_four\_imp\_zero}, through the reflection of
Section~\ref{sec:color}). The gadget accepts an unsatisfying input, and the
judge returns \lean{none}. This is the special case of the paper's Lemma~11
that the minimality proof needs.

\subsection{The search and its completeness}\label{sec:search}
\lean{searchSlice threeColoring parts true} enumerates the symmetric
candidates of a composition. \lean{splitForced} divides the edge orbits
(\lean{Universe.orbits}) into the \emph{base}, the orbits that contain a
forced edge, and the \emph{free} orbits. \lean{searchAux} walks all subsets
of the free orbits depth first, adds each chosen orbit to the base, judges
each leaf with the four-input table, and collects the valid leaves in
\lean{found}. While at least \lean{pruneDepth}${}=8$ free orbits remain
undecided, it skips a subtree when some satisfiable input is already
uncolorable on the partial graph.

\begin{fact}[pruning loses nothing]
\lean{colorsOk} is antitone in the edge list (\lean{colorsOk\_antitone},
from \lean{colorableGoK\_antitone}): removing edges never destroys a proper
coloring. A completion only adds edges, so a satisfiable row that has died
stays dead, and a graph with a dead satisfiable row is rejected
(\lean{commonK\_none\_3col}). Hence \lean{threeColoring\_kill} discharges the
hypothesis of \lean{searchAux\_complete}: every subset of free orbits whose
graph passes the judge is reported.
\end{fact}

\begin{fact}[the walk enumerates the universe's candidates]
The leaf graph of a branch (\lean{walkGraph}) is a permutation of the list
\lean{graphAt base free mask} that \lean{Universe.graphs} produces
(\lean{walkGraph\_perm\_graphAt}), and the judge depends on the edge list
only up to permutation (\lean{checkGadget\_perm\_3col}). So an empty
\lean{found} list means that the judge rejects every graph in
\lean{Universe.graphs}
(\lean{searchSlice\_exhaustive\_threeColoring\_graphs}).
\end{fact}

\subsection{Four inputs suffice for symmetric candidates}\label{sec:symred}
\lean{gadgetK\_s3\_symmetry\_reduction\_threeColoring}
(\lean{ProofsSym.lean}) proves that, on every symmetric candidate, the
four-input judge and the eight-input judge agree. Let $s$ act on inputs by
permuting bits (\lean{envAct}). For an invariant graph, the row of
$s\cdot\mathit{env}$ is the image of the row of $\mathit{env}$ under the
graph automorphism $v\mapsto s\cdot v$: the action maps the edge list onto
itself (\lean{map\_actEdge\_mem}), the vertex list onto itself, and the pins
of $\mathit{env}$ onto those of $s\cdot\mathit{env}$ (the anchors are fixed,
the literals are permuted with their truth values; \lean{rowPins\_colorAt}).
Renaming a 3-coloring instance by an injective map does not change its answer
(\lean{colorable\_rename}). Hence \lean{colorsOk\_row\_act}: the two rows have
the same answer. Every input is $s\cdot r$ for a representative
$r\in\{0,1,3,7\}$ (\lean{exists\_class\_rep}), and the judge's verdict
depends on its rows only through this correspondence
(\lean{commonK\_eq\_of\_quotient}). This is the reduction of the paper's
Section~IV-D that tests one input per class.

\subsection{Every symmetric candidate is enumerated}\label{sec:cover}
\lean{exists\_graph\_of\_invariant} (\lean{ProofsCover.lean}): every
invariant graph inside the edge list that contains the forced edges has the
same edge set as some graph of \lean{Universe.graphs}. Every edge lies in an
emitted orbit (\lean{orbits\_cover}), the action is transitive inside an
orbit (\lean{mem\_orbit\_trans}), so an invariant graph is a union of
orbits; the orbits containing a forced edge are in the base, and the others
are selected by some mask (\lean{exists\_bitSubset}). The judge depends on a
loop-free graph only through its edge set (\lean{gadgetK\_congr\_3col}), so
\lean{searchSlice\_symmetric\_covers\_invariant} follows: if the sweep of a
composition is empty, the eight-input judge rejects every symmetric
candidate of that composition.

\subsection{The trident slice}\label{sec:slice}
\lean{trident\_slice\_found} states that the sweep of \lean{[3, 3]} finds
exactly \lean{[(tridentA, 1), (tridentB, 1)]}. With the facts above,
\lean{trident\_unique\_threeColoring\_invariant} gives: if the judge returns
\lean{some b} on a symmetric candidate of \lean{[3, 3]}, then $b=1$ and its
edge set is that of \lean{tridentA} or \lean{tridentB}.

\subsection{Statements 1--3}
Statements~1 and~2 are closed by \lean{decide}, which the kernel evaluates
(standard axioms only). Statement~3 is \lean{no\_symmetric\_gadget\_below\_six}:
a valid candidate with $k\le5$ would have \lean{parts = [3, 3]} by Statement~4,
whose sizes sum to $6$, a contradiction (\lean{compositions\_sound}).

\subsection{Correspondence with the paper}
The formal proof does not follow the lemma chain of the paper's Section~IV.
Its two proved reductions are general (Sections~\ref{sec:small}
and~\ref{sec:symred}); everything else is settled by exhaustive search,
whose correctness is proved (Sections~\ref{sec:search} and~\ref{sec:cover})
and whose runs are re-done at build time. The paper's own enumeration
(Section~IV-D, Table~1) reports the same 24 compositions.

% ---------------------------------------------------------------------------
\section{The coloring test is correct}\label{sec:color}
Every verdict rests on \lean{colorable}, a backtracking search
(\lean{colorableGo}): it takes the vertices in order, tries the pinned color
or each of $0,1,2$, rejects a color used by an already colored neighbor, and
recurses. \lean{colorable\_iff\_spec} (\lean{ProofsColor.lean}) proves that it
returns \lean{true} exactly when a pin-respecting proper coloring exists
(\lean{SpecExt}), for loop-free graphs and duplicate-free vertex lists (both
hold for candidates: \lean{edges\_ne}, \lean{vertices\_nodup}). Soundness
reads the coloring off a successful run; completeness shows by induction on
the remaining vertices that a run following a given proper coloring is never
rejected. Both are proved once for $k$ colors (\lean{colorableGoK\_sound},
\lean{colorableGoK\_complete}) and instantiated at $k=3$.

% ---------------------------------------------------------------------------
\section{Trusted base}\label{sec:trust}
\begin{itemize}
\item \textbf{Lean's kernel} and the three standard axioms \lean{propext},
\lean{Classical.choice}, \lean{Quot.sound}.
\item \textbf{Mathlib} (v4.33.0) for list lemmas, \lean{interval\_cases},
\lean{omega}, and the group-action library used by Statement~6.
\item \textbf{The compiler, for seven re-runs.} \lean{native\_decide}
evaluates a Boolean by compiled code; Lean records each use as an auxiliary
axiom \lean{<theorem>.\_native.native\_decide.ax\_1\_1}. The seven are
\lean{sym\_empty\_13}, \lean{\_23}, \lean{\_113}, \lean{\_6}, \lean{\_123},
\lean{\_1113} (2\textsuperscript{12} to 2\textsuperscript{26} candidates) and
\lean{trident\_slice\_found} (2\textsuperscript{16}); the kernel evaluator is
impractical beyond about 2\textsuperscript{13} candidates on this engine.
Statements~3 and~4 depend on them; Statements~1, 2 and 5--7 do not.
\end{itemize}
% Axioms printed by `lake build` (the #print axioms lines at the end of
% atlas/Theorem16.lean and atlas/Statement.lean), build of 2026-11-27 on
% Lean v4.33.0.
\medskip
\noindent The build prints the axioms of each result:

\smallskip
\begin{center}
\begin{tabular}{@{}lll@{}}
\toprule
result & standard axioms & compiled re-runs\\
\midrule
\lean{minimum\_unique} (Theorem~16) & all three & all seven\\
\lean{trident\_is\_valid\_gadget} & all three & none\\
\midrule
Statement 1 (the trident is a candidate) & all three & none\\
Statement 2 (the trident is valid) & \lean{propext} & none\\
Statement 3 (minimality) & all three & all seven\\
Statement 4 (uniqueness) & all three & all seven\\
Statement 5 (compositions) & \lean{propext}, \lean{Quot.sound} & none\\
Statement 6 (the action) & all three & none\\
Statement 7 (the judge) & all three & none\\
\bottomrule
\end{tabular}
\end{center}
\smallskip
\noindent ``All three'' are \lean{propext}, \lean{Classical.choice} and
\lean{Quot.sound}. ``All seven'' are the auxiliary axioms
\lean{Atlas.sym\_empty\_13.\_native.native\_decide.ax\_1\_1} and the
corresponding axioms of \lean{sym\_empty\_23}, \lean{sym\_empty\_113},
\lean{sym\_empty\_6}, \lean{sym\_empty\_123}, \lean{sym\_empty\_1113} and
\lean{trident\_slice\_found}. No result depends on \lean{sorryAx}.


% ---------------------------------------------------------------------------
\section{What is not formalized}\label{sec:gaps}
\begin{enumerate}
\item \textbf{From Definition~1 to the candidate space (Lemma~6).} The
statements speak about candidates whose auxiliary vertices come in parts with
the standard action of their size. That every $S_3$-symmetric gadget of
Definition~1 is isomorphic, by a map fixing literals and anchors, to such a
candidate follows from the paper's Lemma~6 (orbit structure) by relabeling
each orbit. The mathematics is standard, but it is not proved in Lean.
Statement~6 shows only that the model's action is an $S_3$ action of the
required kind.
\item \textbf{Uniqueness up to isomorphism.} Statement~4 gives the edge set
of \lean{tridentA} or \lean{tridentB}. That these two are the same graph up
to relabeling the parts, and that the isomorphism of Theorem~16 is the one
fixing literals and anchors, is read off the definitions, not proved.
\item \textbf{The search counts.} The numbers of candidates and leaves in
the paper's Table~1 come from the separate enumeration program; the package
does not prove them.
\end{enumerate}

% ---------------------------------------------------------------------------
\section{Checks}\label{sec:checks}
% Results of the checks, 2026-11-27 (Lean v4.33.0, Mathlib v4.33.0).
\paragraph{Build.} \lean{lake build} completes with 1108 jobs, no errors and
no warnings, and prints the axioms of Section~\ref{sec:trust}. A full build
of the package (with Mathlib taken from its cache) took 8~h~23~min on one
core; nearly all of it is the six compiled re-runs of
\lean{ProofsSweeps16.lean}. Later builds reuse the compiled
result while the sources it depends on are unchanged.

\paragraph{Dependency check.} \lean{checks/ProofCone.lean} follows types,
definition bodies and proofs from the seven statements (core Lean and
Mathlib are not followed) and lists every declaration of the package that is
never reached. Run in the built package with
\lean{lake env lean checks/ProofCone.lean}, it prints
\begin{code}
atlas declarations: 445; used by the seven statements: 346
not used, written in the source: 0
not used, generated by Lean: 99
\end{code}
The 99 unreached declarations are recursors and other constructions that
Lean generates for every structure and inductive type, the derived
\lean{Repr}, \lean{Inhabited} and \lean{DecidableEq} instances, the two
projections \lean{Problem.name} and \lean{Problem.layer} of fields the
statements never read, and one generated congruence lemma. Every
declaration written in the source is used.

\paragraph{Statement cone.} The same traversal restricted to types and
definition bodies gives the declarations whose meaning a reader has to
accept: 56 declarations (257 lines) for Statements~1--4 and 76
declarations (349 lines) for all seven. The README lists their line ranges.


\appendix
\section{Files}
% Files of atlas/ with line counts (2026-11-27).
\begin{center}
\small\begin{tabular}{@{}lrp{0.5\textwidth}@{}}
\toprule
file & lines & content\\
\midrule
\lean{Definition1.lean} & 87 & Definition~1 and the trident graph (Section~\ref{sec:statement}; Mathlib only)\\
\lean{Theorem16.lean} & 257 & \lean{minimum\_unique} and \lean{trident\_is\_valid\_gadget} (Section~\ref{sec:conclusion})\\
\lean{TridentGadget.lean} & 135 & the trident is a valid gadget (Section~\ref{sec:tridentvalid}; Mathlib only)\\
\lean{OrbitStructure.lean} & 173 & Lemma~6 (Section~\ref{sec:lemma6}; Mathlib only)\\
\lean{OrbitLabels.lean} & 177 & permutations and the engine's tables, labels (Section~\ref{sec:labels})\\
\lean{Candidate.lean} & 389 & the orbit composition and the vertex map of a gadget\\
\lean{CandidateGraph.lean} & 156 & the edge list of a gadget is a symmetric candidate\\
\lean{CandidateValid.lean} & 179 & validity is the verdict of the judge (Section~\ref{sec:valid})\\
\midrule
\lean{Statement.lean} & 137 & the seven statements about the candidate space\\
\lean{Meaning.lean} & 155 & proofs of Statements 5--7\\
\lean{Atlas.lean} & 174 & permutations, vertices, universes, the action, edges, orbits, compositions\\
\lean{Framework.lean} & 197 & instances, the problem interface, the judge \lean{gadgetK}, the pruned search\\
\lean{Kits.lean} & 37 & the feasibility kit\\
\lean{Problems/KColoring.lean} & 28 & the backtracking coloring kernel for $k$ colors\\
\lean{Problems/ThreeColoring.lean} & 60 & 3-coloring as a problem: the test and the palette\\
\lean{Proofs.lean} & 168 & the composition enumeration is sound and complete\\
\lean{ProofsLib.lean} & 37 & list lemmas\\
\lean{ProofsAction.lean} & 698 & the action laws, the edge algebra, the orbit generator\\
\lean{ProofsOrbit.lean} & 168 & S3 as a Mathlib group, the \lean{MulAction}s, orbit sizes\\
\lean{ProofsColorK.lean} & 323 & the $k$-coloring kernel: reflection and antitonicity\\
\lean{ProofsColor.lean} & 166 & \lean{colorable\_iff\_spec} (Section~\ref{sec:color})\\
\lean{ProofsPrune.lean} & 44 & pruning soundness\\
\lean{ProofsKits.lean} & 56 & the mask enumeration \lean{bitSubset}\\
\lean{ProofsWalk.lean} & 239 & the search walk and its completeness (Section~\ref{sec:search})\\
\lean{ProofsWalkBridge.lean} & 271 & the walk enumerates the universe's graphs\\
\lean{ProofsSym.lean} & 1123 & four inputs suffice (Section~\ref{sec:symred})\\
\lean{ProofsSmallOrbits.lean} & 257 & orbits of sizes 1 and 2 (Section~\ref{sec:small})\\
\lean{ProofsCover.lean} & 205 & every symmetric candidate is enumerated (Section~\ref{sec:cover})\\
\lean{ProofsDecide.lean} & 160 & \lean{sym\_empty\_3}, the trident slice, \lean{tridentA}, \lean{tridentB}\\
\lean{ProofsSweeps16.lean} & 42 & the six compiled empty sweeps\\
\lean{ProofsTheorem16.lean} & 132 & \lean{trident\_minimum\_unique}\\
\midrule
total & 6430 & \\
\bottomrule
\end{tabular}
\end{center}
The first eight files prove the results for Definition~1; the others are the
candidate-space part. Outside \lean{atlas/}: \lean{checks/ProofCone.lean} (the
dependency check of Section~\ref{sec:checks}) and this report (\lean{report/}).


\end{document}
